\begin{figure}[!t]
\centering
\begin{tikzpicture}[
  box/.style={draw,thin,align=center,font=\small,
    text width=4.15cm,inner sep=5pt},
  wide/.style={box,text width=10.5cm},
  flow/.style={-{Stealth[length=2mm]},thin}]
\node[wide] (endpoint) at (0,0)
  {Assumption~\ref{hyp:mmp}, retaining Assumption~\ref{hyp:campana}\\
   Original nef endpoint \((Y,\Delta)\), \(J=K_Y+\Delta\)};
\node[box,anchor=north] (projective) at (-5,-1.9)
  {Projective case\\
   Lemma~\ref{red:log-iitaka} and Theorem~\ref{proj:thm:main}\\
   Corollary~\ref{red:projective}};
\node[box,anchor=north] (direct) at (0,-1.9)
  {Uniruled or irregular case\\
   Lemmas~\ref{red:uniruled} and~\ref{red:irregular}};
\node[box,anchor=north] (canonical) at (5,-1.9)
  {Remaining smooth canonical case\\
   \(W\) non-uniruled, non-Moishezon, \(q(W)=0\)\par\smallskip
   \(a(W)>0\): Proposition~\ref{pos:nonvanishing}\\
   \(a(W)=0\): Theorem~\ref{end:canonical}, using Sections~\ref{met:section}--\ref{floor:sec}};
\node[wide] (nonvanishing) at (0,-6.3)
  {Theorem~\ref{thm:nonvanishing}: \(\kappa(Y,J)\geq0\)\\
   Nonvanishing on the original nef pair};
\node[wide] (good) at (0,-8.1)
  {Theorem~\ref{thm:main}: a Cartier multiple of \(J\)
   is globally generated on \(Y\)};
\draw[flow] (endpoint.south) -- (0,-1.05) -- (-5,-1.05) -- (projective.north);
\draw[flow] (0,-1.05) -- (direct.north);
\draw[flow] (0,-1.05) -- (5,-1.05) -- (canonical.north);
\node[fill=white,inner sep=2pt,font=\small] at (0,-1.05)
  {Section~\ref{red:section}; Proposition~\ref{red:reduction}};
\draw[flow] (projective.south) -- (-5,-5.6) -- (nonvanishing.north west);
\draw[flow] (direct.south) -- (nonvanishing.north);
\draw[flow] (canonical.south) -- (5,-5.6) -- (nonvanishing.north east);
\draw[flow] (nonvanishing.south) -- node[right,font=\small]
  {Final use of Assumption~\ref{hyp:effective-abundance}} (good.north);
\end{tikzpicture}
\caption{Reading map under Assumptions~\ref{hyp:campana}--\ref{hyp:effective-abundance}.
Lemma~\ref{red:log-iitaka} derives the sole premise of Theorem~\ref{proj:thm:main}
from Assumption~\ref{hyp:campana},
so the projective appendices add no independent conditional input.
The groupings indicate uses in this proof, not the full scope of the
component lemmas. Assumption~\ref{hyp:effective-abundance} also enters
intermediate constructions. All routes return nonvanishing to the
original nef pair before its final application on that endpoint.}
\label{fig:proof-map}
\end{figure}
